En faisant réagir $\qty{11.2}{\g}$ de fer avec un volume $V=\qty{6.0}{\L}$ de
dichlore, on obtient du chlorure de fer(III) solide. Dans les conditions de
l'expérience, le volume molaire vaut $V_{m}=\qty{24}{\L\per\mol}$.
\begin{enumerate}
    \item Écrire l'équation de la réaction.
    \item Calculer les quantités initiales des réactifs.
    \item Quel est le réactif limitant~?
\end{enumerate}
